\begin{theorem}[结论 10（2026-04-14 12:43:32 +08）：共振缺口的不可约起源与最小可见模块的单纯性]\label{thm:killo-fold-resonance-gap-simple-module}
\leanverified{paper\_killo\_fold\_resonance\_gap\_simple\_module}
对 \(q\in\{13,14,15,16,17\}\) 记
\[
N_q:=2\lfloor q/2\rfloor+1,\qquad
d_q:=\deg P_q,\qquad
\Delta(q):=N_q-d_q.
\]
则
\[
\Delta(13)=2,\quad \Delta(14)=2,\quad \Delta(15)=4,\quad \Delta(16)=4,\quad \Delta(17)=4.
\]
并且任取序列 \(m\mapsto S_q(m)\) 的一个最小有理线性 realization
\[
S_q(m)=u_q^{\mathsf T}A_q^{\,m}v_q,\qquad A_q\in M_{d_q}(\QQ),
\]
则 \(\QQ^{d_q}\) 作为 \(\QQ[A_q]\)-模同构于域
\[
\QQ[\lambda]/(P_q(\lambda)),
\]
从而是一个单纯 \(\QQ[A_q]\)-模。
特别地，\(A_q\) 不存在任何非平凡的有理不变直和分解
\[
A_q\sim B\oplus C,\qquad 0<\dim B,\dim C<d_q.
\]
因此，正缺口 \(\Delta(q)>0\) 不能由最小可见特征式的有理因式分解或可见状态块的有理直和拆分来解释；它必然发生在最小特征式形成之前的可见函子压缩阶段。
\end{theorem}

\begin{proof}
由表 \ref{tab:fold_collision_resonance_gap_delta_q_9_17} 直接读出
\[
(d_{13},d_{14},d_{15},d_{16},d_{17})=(11,13,11,13,13),
\]
而
\[
(N_{13},N_{14},N_{15},N_{16},N_{17})=(13,15,15,17,17),
\]
故缺口数值如上。

再由命题 \ref{prop:pom-moment-minrecurrence-hankel} 与定理
\ref{thm:hankel-rank-minimal}，
最小 realization 的维数正等于最小递推阶数 \(d_q\)，且其最小多项式就是 \(P_q\)。
另一方面，命题 \ref{prop:pom-resonance-window-minpoly-irreducible-q9-17} 给出
\(
P_q
\)
在 \(\QQ\) 上不可约。
由于
\(
\dim \QQ^{d_q}=d_q=\deg P_q
\)，
\(A_q\) 的有理标准形只能由单个不变因子 \(P_q\) 构成；亦即
\[
\QQ^{d_q}\cong \QQ[\lambda]/(P_q(\lambda))
\]
作为 \(\QQ[\lambda]\)-模同构。
而 \((P_q)\subset \QQ[\lambda]\) 为极大理想，故商环
\(
\QQ[\lambda]/(P_q)
\)
是域；域作为其自身上的模没有非平凡子模。
于是 \(\QQ^{d_q}\) 是单纯 \(\QQ[A_q]\)-模。

若存在非平凡有理不变直和分解 \(A_q\sim B\oplus C\)，则
\(
\QQ^{d_q}
\)
将分解为两个非零 \(\QQ[A_q]\)-子模，与单纯性矛盾。
因此最小可见动力学内部不存在任何有理块分裂。

最后，若正缺口 \(\Delta(q)>0\) 可由最小可见特征式层面的有理因式分解或可见状态块删减产生，则在最小 \(d_q\)-维 realization 上必残留一个非平凡有理不变块或非平凡首一因子；这与上面的单纯性及 \(P_q\) 的不可约性矛盾。
故缺口只能发生在最小特征式形成之前，即发生于可见函子对名义模态空间的压缩阶段。
证毕。
\end{proof}

\begin{corollary}[共振缺口不是伽罗瓦退化]\label{cor:killo-fold-resonance-gap-not-galois-degeneracy}
\leanverified{paper\_killo\_fold\_resonance\_gap\_not\_galois\_degeneracy}
对 \(q\in\{13,14,15,16,17\}\)，正缺口 \(\Delta(q)>0\) 与
\[
\Gal(P_q/\QQ)\cong S_{d_q}
\]
同时成立。
因而该缺口并非来自可解子扇区、判别式平方化、或可见特征式的算术可分性，而是来自可见态函子在形成最小谱对象之前已经执行的严格压缩。
\end{corollary}

\begin{proof}
对 \(q=13,14,15\) 由定理 \ref{thm:pom-resonance-minpoly-galois-sd-q12-15}，
对 \(q=16,17\) 由定理 \ref{thm:pom-resonance-galois-s13-q16-17}，皆有满对称伽罗瓦群。
与定理 \ref{thm:killo-fold-resonance-gap-simple-module} 合并即得结论。
\end{proof}
